\section{Bias-Driven Dynamic-Range Inflation in Block Quantization}
\label{app:quant_bias}

This appendix formalizes how activation bias inflates blockwise quantization
scales and thereby degrades the effective resolution available to long-tail
components.

\subsection{Setup: block quantization model}

Consider a vector $x\in\mathbb{R}^H$ that will be quantized using blockwise
scaling.
Let $B\subset\{1,\dots,H\}$ denote a block of coordinates with $|B|=b$.
A standard block quantizer chooses a scale proportional to
\begin{equation}
\label{eq:block_scale_max}
\alpha(x;B) \triangleq \max_{j\in B} |x_j|.
\end{equation}
The quantization step size is then
\[
\Delta(x;B) = \alpha(x;B) / Q,
\]
where $Q$ is the maximum representable integer level (e.g.\ $Q=7$ for signed
4-bit symmetric quantization).

We study how $\alpha(x;B)$ behaves under the decomposition
\begin{equation}
\label{eq:x_bias_noise}
x_j = \mu_j + s_j,
\qquad
\mu\in\mathbb{R}^H \text{ deterministic},
\qquad
\mathbb{E}[s_j]=0.
\end{equation}

\subsection{Assumptions}

\begin{assumption}[Sub-Gaussian centered variation]
\label{assump:subg_C}
For each $j$, the centered noise $s_j$ is sub-Gaussian with parameter $\sigma^2$:
\[
\mathbb{P}(|s_j|\ge t)\le 2\exp\!\left(-\frac{t^2}{2\sigma^2}\right).
\]
The $s_j$ are independent across coordinates within a block.
\end{assumption}

The independence assumption is used only to obtain clean max bounds; weaker
dependence conditions yield similar results with different constants.

\subsection{Block scale with bias}

Define
\[
m_B \triangleq \max_{j\in B} |\mu_j|.
\]

\begin{lemma}[Upper bound]
\label{lem:block_upper}
For any $t\ge 0$, with probability at least $1-2b\exp(-t^2/(2\sigma^2))$,
\[
\alpha(x;B) \le m_B + t.
\]
\end{lemma}

\begin{proof}
For any $j\in B$, $|x_j| \le |\mu_j| + |s_j|$.
Hence
\[
\max_{j\in B}|x_j|
\le
\max_{j\in B}|\mu_j| + \max_{j\in B}|s_j|
=
m_B + \max_{j\in B}|s_j|.
\]
By a union bound,
\[
\mathbb{P}\left(\max_{j\in B}|s_j|\ge t\right)
\le
\sum_{j\in B}\mathbb{P}(|s_j|\ge t)
\le
2b\exp\!\left(-\frac{t^2}{2\sigma^2}\right).
\]
\end{proof}

\begin{lemma}[Lower bound]
\label{lem:block_lower}
For any $t\ge 0$, with probability at least $1-2\exp(-t^2/(2\sigma^2))$,
\[
\alpha(x;B) \ge m_B - t.
\]
\end{lemma}

\begin{proof}
Let $j^\star\in B$ achieve $|\mu_{j^\star}|=m_B$.
Then
\[
\alpha(x;B)=\max_{j\in B}|x_j|\ge |x_{j^\star}|
\ge |\mu_{j^\star}| - |s_{j^\star}|
= m_B - |s_{j^\star}|.
\]
Thus
\[
\mathbb{P}\big(\alpha(x;B)\le m_B - t\big)
\le
\mathbb{P}(|s_{j^\star}|\ge t)
\le
2\exp\!\left(-\frac{t^2}{2\sigma^2}\right).
\]
\end{proof}

\paragraph{Consequence.}
Combining Lemmas~\ref{lem:block_upper} and \ref{lem:block_lower} shows that
block scale concentrates around $m_B$ whenever $m_B \gg \sigma\sqrt{\log b}$.

\subsection{Block scale without bias (mean removal)}

After mean removal, the centered vector is $x'_j = s_j$ and the block scale is
\[
\alpha(x';B) = \max_{j\in B}|s_j|.
\]

\begin{lemma}[Centered max bound]
\label{lem:centered_max}
For any $\delta\in(0,1)$, with probability at least $1-\delta$,
\[
\alpha(x';B) \le \sigma \sqrt{2\log(2b/\delta)}.
\]
\end{lemma}

\begin{proof}
Set $t=\sigma\sqrt{2\log(2b/\delta)}$ in Lemma~\ref{lem:block_upper} with $m_B=0$:
\[
\mathbb{P}\!\left(\max_{j\in B}|s_j|\ge t\right)
\le 2b\exp\!\left(-\frac{t^2}{2\sigma^2}\right)
= 2b\exp\!\left(-\log(2b/\delta)\right)=\delta.
\]
\end{proof}

\subsection{Dynamic-range inflation factor}

The ratio of block scales with and without bias is
\[
\frac{\alpha(x;B)}{\alpha(x';B)}.
\]

\begin{theorem}[Bias-driven dynamic-range inflation]
\label{thm:range_inflation}
Fix $\delta\in(0,1)$.
Under Assumption~\ref{assump:subg_C}, with probability at least $1-2\delta$,
\[
\frac{\alpha(x;B)}{\alpha(x';B)}
\ge
\frac{m_B - \sigma\sqrt{2\log(2/\delta)}}{\sigma\sqrt{2\log(2b/\delta)}},
\]
and
\[
\frac{\alpha(x;B)}{\alpha(x';B)}
\le
\frac{m_B + \sigma\sqrt{2\log(2b/\delta)}}{\sigma\sqrt{2\log(2b/\delta)}}.
\]
In particular, if $m_B \gg \sigma\sqrt{\log b}$, then
\[
\alpha(x;B) \approx m_B
\qquad\text{while}\qquad
\alpha(x';B) = O(\sigma\sqrt{\log b}),
\]
so the block scale is inflated by a factor on the order of $m_B / (\sigma\sqrt{\log b})$.
\end{theorem}

\begin{proof}
The upper bound follows by combining Lemma~\ref{lem:block_upper} (with probability
$1-\delta$) and Lemma~\ref{lem:centered_max} (with probability $1-\delta$).
The lower bound follows by Lemma~\ref{lem:block_lower} and Lemma~\ref{lem:centered_max}.
A union bound gives total probability at least $1-2\delta$.
\end{proof}

\subsection{Implication for quantization resolution and tail SNR}

Block quantization step size satisfies
\[
\Delta(x;B) = \alpha(x;B)/Q.
\]
Thus Theorem~\ref{thm:range_inflation} implies that bias increases the quantization
step size multiplicatively in blocks containing large-mean coordinates.
This reduces effective resolution for coordinates whose magnitude is governed by
centered variation $s_j$, i.e., the long-tail semantic component.

A crude but informative proxy for tail signal-to-noise ratio (SNR) is
\[
\mathrm{SNR}_{\text{tail}} \propto \frac{\mathbb{E}[s_j^2]}{\Delta(x;B)^2},
\]
which decreases quadratically with block scale.
Therefore, bias-driven inflation of $\alpha(x;B)$ directly suppresses tail SNR.

\paragraph{Remark (percentile-based scaling).}
In practice, scales are often set by a high percentile (e.g.\ $p99.9$) rather than
the maximum.
The same phenomenon holds: when the mean dominates the upper quantiles of $|x_j|$,
percentile-based scales concentrate around $m_B$ rather than the noise level.
Replacing the max bounds above with sub-Gaussian quantile bounds yields analogous
conclusions.

\paragraph{Summary.}
Bias inflates blockwise quantization scales whenever blocks contain coordinates
with large $|\mu_j|$, compressing the resolution allocated to tail variability.
Mean removal eliminates this inflation and restores block scales to the
noise-controlled regime $O(\sigma\sqrt{\log b})$.
This formalizes the empirical dynamic-range reductions observed in
Section~2.4C.